\chapter{Grothendieck local duality}\label{ch:local-duality}

We use cohomological indexing, so $K[a]^j=K^{j+a}$ and
$d_{K[a]}=(-1)^a d_K$. A module without an indicated shift is
in degree zero. Write $\mathcal H om(X,Y)$ for the Hom complex
\[
 \mathcal H om(X,Y)^a=\prod_j\operatorname{Hom}(X^j,Y^{j+a}),
 \qquad d(f)=d_Y f-(-1)^a f d_X .
\]
Throughout, $C_{\mathfrak m}$ denotes the augmented \v{C}ech
complex of the local ring in question, with the ring in degree
zero, as in \ref{res:L05}. We write $\operatorname{Ext}^j=0$ for
negative degrees whenever both arguments are modules. References
for the local duality formulas are \cite[Sections 18--22]{st-dual}
and \cite[Section 3.5]{bruns}.

\section{Duality over a Gorenstein local ring}

\begin{proposition}[The top local cohomology of a Gorenstein local ring]\label{res:D01}
Let $(S,\mathfrak n,k)$ be Gorenstein of dimension $n$. Then
$R\Gamma_{\mathfrak n}(S)\simeq E_S(k)[-n]$.
\end{proposition}
\begin{proof}
For this ring, $\operatorname{Ext}^i_S(k,S)$ is zero for $i\ne n$
and is a one-dimensional $k$-space for $i=n$
\cite[GC-06]{core}. In the minimal injective resolution from
\ref{res:M07}, these Ext groups are its socles. Thus applying
$\Gamma_{\mathfrak n}$ gives a complex with a single nonzero term,
a copy of $E_S(k)$ in degree $n$. This complex computes the derived
functor by the injective resolution. Choosing a different
identification with the chosen injective hull postcomposes this
comparison by an automorphism of that hull.
\end{proof}

\begin{lemma}[A finite-free model for the dual calculation]\label{res:D02}
Under the hypotheses of \ref{res:D01}, let
$P\to X$ be a bounded-above, degreewise finite free resolution of
a finite $S$-module $X$, and put $E_S=E_S(k)$. There are
quasi-isomorphism comparisons
\[
 P\otimes_S C_{\mathfrak n}\simeq X\otimes_S C_{\mathfrak n},
 \qquad
 P\otimes_S C_{\mathfrak n}\simeq P\otimes_S E_S[-n].
\]
After applying the dual into $E_S[0]$, the last model is identified
by tensor--Hom adjunction with $\mathcal H om_S(P,\widehat S[n])$.
\end{lemma}
\begin{proof}
Choose $P$ by successively resolving finite kernels. A bounded
complex of flat modules preserves acyclicity under tensor product:
filter by the subcomplexes obtained by retaining only terms
in degrees greater than or equal to a given index. Each successive
quotient is one flat module in one degree, and successive short
exact sequences are split in each degree. The same argument applies
to a bounded-above complex via the increasing union of its bounded
degree truncations. Indeed its tensor product is the filtered
colimit of the tensor products of the truncations, and filtered
colimits of modules are exact. This proves both applications of
tensoring a quasi-isomorphism, the first using the bounded flat
$C_{\mathfrak n}$ and the second using $P$ with \ref{res:D01}.
One can use quasi-isomorphisms in either direction through a
common complex, since tensoring preserves each of them.

In the model with $E_S[-n]$, each total degree in the tensor
product uses a single term of $P$. With the usual tensor
differential $d(p\otimes e)=dp\otimes e+(-1)^{\deg p}p\otimes de$,
the assignment $f\mapsto (p\mapsto(e\mapsto f(p\otimes e)))$
is a chain isomorphism
\[
 \mathcal H om_S(P\otimes_S E_S[-n],E_S[0])
 \ \cong\
 \mathcal H om_S(P,\mathcal H om_S(E_S[-n],E_S[0])) .
\]
In the differential, the $d_P$ term has sign $-(-1)^{\deg f}$,
and the $d_{E_S[-n]}$ term has sign
$-(-1)^{\deg f+\deg p}$ on both sides. The inner Hom on the right
is concentrated in degree $-n$ and has value
$\operatorname{End}_S(E_S)=\widehat S$ by \ref{res:M03}. This
gives the indicated shift.
\end{proof}

\begin{theorem}[Gorenstein local duality]\label{res:D03}
Under the assumptions of \ref{res:D01}, for each finite
$S$-module $X$ there are natural isomorphisms
\[
 \bigl(H^i_{\mathfrak n}(X)\bigr)^\vee
   \ \cong\ \operatorname{Ext}^{n-i}_S(X,\widehat S)
   \ \cong\ \widehat{\operatorname{Ext}^{n-i}_S(X,S)} .
\]
\end{theorem}
\begin{proof}
Dualizing into the injective module $E_S$ is exact. On complexes it gives
\[
     H^j(\mathcal H om_S(U,E_S[0]))\cong H^{-j}(U)^\vee.
\]
For example, a cocycle of degree $j$ vanishes on boundaries in
degree $-j$ and a coboundary vanishes on the cycles. Every map on
those cycle classes extends to the ambient quotient by injectivity.
This proves the identification and its naturality. Now apply it to
\ref{res:L05} and take degree $-i$ in the comparison of
\ref{res:D02}. In the last complex, shifting the second argument
by $n$ shifts the Hom complex by $n$, so the result is
$\operatorname{Ext}^{n-i}_S(X,\widehat S)$.

Completion is flat and agrees with tensoring by $\widehat S$ on
finite modules \cite[ALG-05-C04-R03 and S01]{p103}. For each term
of the degreewise finite free resolution, base change of its
$S$-dual gives its $\widehat S$-valued dual. Flatness commutes with
taking cohomology, giving the last isomorphism. A map of finite
modules lifts to a chain map of free resolutions by lifting into
successive cycle modules, starting in degree zero. Any two lifts
are chain homotopic by lifting their difference into successive
boundary preimages. The chain comparisons therefore give
naturality on modules. With the scalar action induced by
\ref{res:M03}, the maps are also $\widehat S$-linear.
\end{proof}

\section{The general complete case}

\begin{lemma}[Adjunction for a local quotient]\label{res:D04}
If $S\twoheadrightarrow R$ is a surjective local map of Noetherian
local rings with residue field $k$, put $\mathfrak m=\mathfrak nR$.
There is an identification
\[
 E_R(k)\cong\operatorname{Hom}_S(R,E_S(k)).
\]
For an $R$-module $X$, the local cohomology over $R$ is computed by
the \v{C}ech model over $S$ on $X$, and
\[
 \operatorname{Hom}_R(X,\operatorname{Hom}_S(R,E_S(k)))
   \ \cong\ \operatorname{Hom}_S(X,E_S(k)).
\]
Both comparisons respect the $R$-action.
\end{lemma}
\begin{proof}
Use precomposition with multiplication for the $R$-action on the
inner Hom. The displayed adjunction is evaluation at $1$ with
inverse $h\mapsto [x\mapsto (r\mapsto h(rx))]$. Hence the inner Hom
is injective over $R$. Evaluating it at $1$ identifies it as the
elements of $E_S(k)$ annihilated by the kernel of $S\to R$.
It contains $k$ essentially, because it is a submodule of the
$S$-injective hull of $k$ and its $S$-submodules are its
$R$-submodules. It is therefore an injective hull over $R$.
If $s_1,\ldots,s_l$ generate $\mathfrak n$, their images generate
$\mathfrak m$. Each term of the \v{C}ech complex on $X$
localizes $X$ at the same endomorphism of $X$ for either choice
of scalar ring, so \ref{res:L08} gives the assertion. All the
maps commute both with multiplication by $R$ on $X$ and with
the adjunction.
\end{proof}

\begin{theorem}[A normalized dualizing complex]\label{res:D05}
Let $R$ be complete Noetherian local, and let $S\twoheadrightarrow R$ be a surjection
from a complete regular local ring of dimension $n$ as in \ref{res:C04}.
Put
\[
          \mathcal D_R=\mathrm{RHom}_S(R,S[n]).
\]
It can be represented by a bounded complex of injective $R$-modules,
and its cohomology modules are finite. The homothety map
$R\longrightarrow\mathrm{RHom}_R(\mathcal D_R,\mathcal D_R)$ is an
isomorphism, and $\mathrm{RHom}_R(k,\mathcal D_R)\simeq k[0]$.
\end{theorem}
\begin{proof}
The regular local ring $S$ is Gorenstein and has injective
dimension $n$ as a module over itself \cite[GC-06--GC-07]{core}.
Represent $S[n]$ by a bounded injective resolution $J$ and let
$I=\mathcal H om_S(R,J)$, with $R$ acting on the argument $R$.
Each term of $I$ is injective by the adjunction of \ref{res:D04}.
A bounded complex of injectives is a valid target for derived
Hom. Filter it by its bounded degree truncations, use degreewise
split short exact sequences of complexes, and use exactness of
Hom into each single injective term to see that Hom from any
acyclic complex into it is acyclic. Thus the termwise adjunction
gives
\[
        \mathrm{RHom}_R(U,\mathcal D_R)
        \ \simeq\ \mathrm{RHom}_S(U,S[n])
\]
for every $R$-complex $U$. At $U=k$ it gives the normalization by
the Gorenstein residue-field Ext calculation \cite[GC-06]{core}.
Finiteness of $H^j(I)$ follows from computing $\mathrm{RHom}_S(R,S[n])$
by a degreewise finite free resolution over the Noetherian ring.
The cohomology is an $R$-module, so finiteness over $S$ implies
finiteness over $R$.

For the homothety assertion, the same adjunction at $U=I$
turns scalar multiplication by $r$ on $I$ into
$[i\mapsto i(r)]\in\mathcal H om_S(I,J)$. This is the natural
derived biduality morphism of $R$ over $S$ with target $S[n]$.
Here is a direct comparison. Since $S$ is regular local, the
finite module $R$ has a bounded finite free resolution
$P\to R$ \cite[CM-12]{core}. Set
$Q=\mathcal H om_S(P,S[n])$. The maps
$I\to\mathcal H om_S(P,J)\leftarrow Q$ are quasi-isomorphisms.
The left one uses that $J$ is a bounded injective target.
The right one follows from filtering the bounded finite free
source $P$ and using exactness of Hom from each free term.
For homogeneous $p\in P^a$ and $q\in Q^b$, the bidual pairing
sends $p$ against $q$ to $(-1)^{ab}q(p)$. Substituting into the
Hom differential verifies it is a chain map. Its map
\[
          P\longrightarrow \mathcal H om_S(Q,S[n])
\]
is an isomorphism, since in each degree it is finite free
biduality up to this sign. The finite free terms of $Q$ give
$\mathcal H om_S(Q,S[n])\simeq \mathcal H om_S(Q,J)$, and taking
Hom into $J$ in the preceding zigzag compares this target to
$\mathcal H om_S(I,J)$. The evaluation pairing is natural in
the morphisms of the source and the target, so under these
comparisons evaluation on $P$ gives $i\mapsto i(r)$ after
$P\to R$. This proves that the actual homothety morphism is
a quasi-isomorphism.
\end{proof}

\begin{theorem}[Local duality for a complete local ring]\label{res:D06}
For every complete Noetherian local ring $R$ and finite $R$-module $X$,
\[
 \bigl(H^i_{\mathfrak m}(X)\bigr)^\vee
     \ \cong\ H^{-i}\!\left(\mathrm{RHom}_R(X,\mathcal D_R)\right)
\]
naturally as $R$-modules. With a chosen $E_R(k)$,
$\mathcal H om_R(C_{\mathfrak m},E_R(k)[0])$ represents
$\mathcal D_R$.
\end{theorem}
\begin{proof}
Choose a complete regular cover $S$ by \ref{res:C04}.
Using the \v{C}ech complex over $S$, consider
\[
                W=\mathcal H om_S(C_{\mathfrak n},E_S[0]).
\]
It is bounded and its terms are injective. For each flat term
$F$ of $C_{\mathfrak n}$, the functor
$\operatorname{Hom}_S(-,\operatorname{Hom}_S(F,E_S))$
is exact by tensor--Hom adjunction, flatness of $F$, and
injectivity of $E_S$. By \ref{res:D01}, \ref{res:D03} with
$X=S$, and completeness, $W$ has just one cohomology module,
namely $S$ in degree $-n$. A cocycle representing a generator
defines a quasi-isomorphism $S[n]\to W$. Thus the bounded
injective complex $W$ may be used as $J$ in \ref{res:D05}.
Choose the generators of $\mathfrak m$ as the images of the
generators of $\mathfrak n$. Termwise adjunction, with the
complex signs of \ref{res:D02}, gives the isomorphism of
$R$-complexes
\[
 \begin{aligned}
 \mathcal H om_S(R,W)
 &\cong \mathcal H om_S(R\otimes_S C_{\mathfrak n},E_S[0])\\
 &\cong \mathcal H om_R(C_{\mathfrak m},E_R(k)[0]).
 \end{aligned}
\]
The second line uses \ref{res:D04}, together with the ordinary
base change of the flat localized terms of the \v{C}ech complex.
In particular the action on the right is the usual $R$-action,
and on the left it is the action on the argument $R$.
Any choice of generators gives the same derived torsion
comparison by \ref{res:L05}, and exact dualization preserves
those quasi-isomorphisms.

The displayed bounded injective complex computes Hom out of $X$
by plain Hom. Under adjunction it is
$\mathcal H om_R(X\otimes_R C_{\mathfrak m},E_R(k)[0])$.
By \ref{res:L05} and the exact-dual cohomology calculation in
\ref{res:D03}, its degree $-i$ cohomology is the asserted
Matlis dual. All identifications after a choice of the dualizing
model are natural in $X$.
\end{proof}

\begin{corollary}[Local duality after completion]\label{res:D07}
If $R$ is any Noetherian local ring and $X$ is finite, then as
$\widehat R$-modules
\[
 \bigl(H^i_{\mathfrak m}(X)\bigr)^\vee
 \ \cong\
 H^{-i}\!\left(\mathrm{RHom}_{\widehat R}
        (\widehat X,\mathcal D_{\widehat R})\right).
\]
\end{corollary}
\begin{proof}
The completion comparison \ref{res:L09} identifies
$H^i_{\mathfrak m}(X)$ with
$H^i_{\mathfrak m\widehat R}(\widehat X)$ with its unique
completed action. Choose the $\widehat R$-injective hull of
$k$ as $R$-injective hull after restriction, as in \ref{res:M06}.
For torsion modules all $R$-linear maps to it are already
$\widehat R$-linear, so this comparison identifies the two
Matlis duals over $\widehat R$. Now use \ref{res:D06} for
$\widehat R$.
\end{proof}

\section{Canonical modules and examples}

\begin{proposition}[Canonical modules and concentration]\label{res:D08}
Let $R$ be a complete Noetherian local ring with $\dim R=d$. Then
\[
      \omega_R=H^{-d}(\mathcal D_R)
             \ \cong\ (H^d_{\mathfrak m}(R))^\vee
\]
is finite, with $\omega_R^\vee\cong H^d_{\mathfrak m}(R)$. In
particular it is the canonical module in the sense of being the
unique finite module up to isomorphism with this dual. The ring
$R$ is Cohen--Macaulay if and only if
$\mathcal D_R\simeq\omega_R[d]$. In that case $\omega_R$ is also
the canonical module of \cite[GC-10]{core}, computed as
$\operatorname{Ext}^{n-d}_S(R,S)$ for a regular cover as above.
\end{proposition}
\begin{proof}
Apply \ref{res:D06} at $X=R$. The stated dual characterization
and its uniqueness follow from \ref{res:M05} and the Artinianness
in \ref{res:M08}. By detection in \ref{res:M01}, nonzero
cohomology of $\mathcal D_R$ corresponds in opposite degrees to
nonzero local cohomology of $R$. By the vanishing and
Cohen--Macaulay criterion \ref{res:V07}, only degree $-d$ is
nonzero exactly in the Cohen--Macaulay case. A complex with a
single nonzero cohomology module is isomorphic in the derived
category to that module in the corresponding degree. To see this,
truncate at that degree with the kernel of the outgoing
differential as the last term, then quotient that kernel by
boundaries to obtain the indicated zigzag. Finally, shifting
the second argument in the formula defining $\mathcal D_R$ in
\ref{res:D05} gives the stated Ext identification. It is the
same module and quotient-ring action as in \cite[GC-10]{core}.
\end{proof}

\begin{corollary}[Cohen--Macaulay duality]\label{res:D09}
For a complete Cohen--Macaulay local ring of dimension $d$
and finite $R$-module $X$ there are natural identifications
\[
 (H^i_{\mathfrak m}(X))^\vee
       \cong \operatorname{Ext}^{d-i}_R(X,\omega_R),\qquad
 H^i_{\mathfrak m}(X)
       \cong (\operatorname{Ext}^{d-i}_R(X,\omega_R))^\vee .
\]
For noncomplete Cohen--Macaulay $R$ the same comparisons are
available using $\widehat R,\widehat X,\omega_{\widehat R}$.
If a finite module $\omega$ over $R$ is given with
$\widehat\omega\simeq\omega_{\widehat R}$, the dual can equivalently
be written as $\widehat{\operatorname{Ext}^{d-i}_R(X,\omega)}$.
\end{corollary}
\begin{proof}
After substituting \ref{res:D08} in \ref{res:D06}, the shift adds $d$ to
the index of the Hom complex. Artinianness \ref{res:M08} and
evaluation \ref{res:M05} give the second formula. In the
noncomplete case dimension and depth of a finite module
are preserved by completion, in particular by \ref{res:V04}
applied to $R$. Thus $\widehat R$ is Cohen--Macaulay of
dimension $d$. Use \ref{res:D07} and Artinian-side evaluation
from \ref{res:M06}. Finally, in a degreewise finite free
resolution of $X$, flat base change gives a resolution of
$\widehat X$. In each term it commutes with the Hom into
$\omega$, hence with cohomology. The Ext groups on the
original ring are finite, so their completed modules are the
corresponding tensors. This proves the asserted variant.
\end{proof}

\begin{example}[Power series and a noncomplete DVR]\label{res:E01}
For $S=k[[x_1,\ldots,x_d]]$, the inverse monomials
$x_1^{-a_1}\cdots x_d^{-a_d}$ with all $a_j\ge1$ form a $k$-basis
of $H^d_{\mathfrak n}(S)\cong E_S(k)$. For $T=k[t]_{(t)}$,
$H^1_{(t)}(T)=k(t)/T$ and its Matlis dual is $k[[t]]$.
\end{example}
\begin{proof}
For $d>0$ the top \v{C}ech quotient is
\[
 S_{x_1\cdots x_d}\bigg/\sum_j S_{\prod_{i\ne j} x_i}.
\]
With denominator $\prod_i x_i^N$, retain only numerator terms
whose exponents in all variables are $<N$. The remaining terms
can be assigned to the first index $j$ at which the exponent is
at least $N$. On each such subseries the denominator in $x_j$
cancels, so its class is zero. The retained terms give a finite
linear combination of inverse monomials. Conversely, after
putting a finite linear combination and a sum of terms from the
smaller localizations over a common denominator, coefficient
comparison in the formal series ring forces the coefficients
of the inverse monomials to vanish. For $d=0$ the empty product
is the basis of $k$. The ring is Gorenstein, so \ref{res:D01}
identifies this top cohomology with $E_S(k)$.
For the DVR $T$, the one-generator complex is
$T\to k(t)$, so its first cohomology is the asserted quotient.
It is $E_T(k)$ by \ref{res:D01}, with dual the completion by
\ref{res:M03}.
\end{proof}

\begin{example}[An Artinian dualizing module]\label{res:E02}
For $R=k[[x,y]]/(x,y)^2$, one has
$\mathcal D_R\simeq\operatorname{Hom}_k(R,k)[0]$.
The ring has socle dimension two, and the dualizing module
has socle dimension one.
\end{example}
\begin{proof}
As a $k$-space $R$ has basis $1,x,y$ and
$\mathfrak m=\langle x,y\rangle$ with square zero.
Thus $\operatorname{soc}R=\mathfrak m$. The $k$-dual has the
usual $R$-action by precomposition and is injective by the
adjunction $\operatorname{Hom}_R(U,\operatorname{Hom}_k(R,k))
\cong\operatorname{Hom}_k(U,k)$. Its socle consists of the
functionals vanishing on $\mathfrak m$, a single copy of $k$.
It is essential, since every module is power torsion for this
nilpotent maximal ideal and \ref{res:M04} applies. Hence it is
$E_R(k)$. The \v{C}ech complex on $x,y$ has only the ring
term nonzero, so \ref{res:D06} computes the asserted dualizing
complex. The Artinian socle criterion
\cite[GC-05--GC-06]{core} shows that $R$ is not Gorenstein.
\end{proof}

\begin{example}[A one-dimensional hypersurface]\label{res:E03}
Let $R=k[[u,v]]/(uv)$. Then $\omega_R\cong R$. For $X=R$,
$H^0_{\mathfrak m}(X)=0$ and $H^1_{\mathfrak m}(X)=E_R(k)$.
For $X=k$, $H^0_{\mathfrak m}(X)=k$ and $H^1_{\mathfrak m}(X)=0$.
\end{example}
\begin{proof}
The element $uv$ is regular in the two-dimensional regular
local ring $k[[u,v]]$, so $R$ is a one-dimensional Gorenstein
local ring \cite[CM-02 and GC-07]{core}. The two-term free
resolution of $R$ over $k[[u,v]]$ has differential $uv$ and its
degree-one Ext into that ring is $R$. This is the canonical
module by \ref{res:D08}. For $X=R$, \ref{res:D01} gives the
cohomology assertions. In coordinates, $R$ consists of pairs
of series in $k[[u]]\oplus k[[v]]$ with equal constant terms.
The \v{C}ech complex has first map the inclusion into
$k((u))\oplus k((v))$ and next term zero since $uv=0$.
The degree-one quotient realizes the hull explicitly. For
$X=k$, all the localizations at either generator vanish, giving
the assertions by \ref{res:L05}. The degrees can also be seen
from \ref{res:D09}. The Ext groups of $R$ into $R$ in degrees
zero and one are $R$ and $0$, respectively, and those of $k$
into $R$ are $0$ and $k$, respectively, by
\cite[GC-06]{core}.
\end{proof}

\begin{example}[A dualizing complex with two nonzero degrees]\label{res:E04}
Let $S=k[[x,y]]$ and $R=S/(x^2,xy)$. Then
\[
 \begin{array}{ll}
 H^0_{\mathfrak m}(R)=kx, &
 H^1_{\mathfrak m}(R)\cong k((y))/k[[y]],\\[2pt]
 H^{-1}(\mathcal D_R)\cong S/(x), &
 H^0(\mathcal D_R)\cong k .
 \end{array}
\]
All remaining cohomology of $\mathcal D_R$ is zero.
\end{example}
\begin{proof}
In $R$, the element $x$ is nilpotent and killed by $y$.
Consequently the \v{C}ech model on $x,y$ reduces to $R\to R_y$.
The ring is the direct sum $k[[y]]\oplus kx$ as a $k$-space,
and localization kills $kx$ and embeds $k[[y]]$ into $k((y))$.
This gives the local cohomology calculation.

Over $S$ a free resolution has consecutive maps
\[
        S\xrightarrow{(-y,x)^{\mathsf T}} S^2
          \xrightarrow{(x^2,xy)} S \longrightarrow R .
\]
The first map is injective. If $x^2a+xyb=0$, cancel $x$, reduce
modulo $x$ to get $b=xc$, and conclude $a=-yc$. Thus the maps
are exact where required. Upon dualizing into $S$, the maps
are the column $(x^2,xy)^{\mathsf T}$ and the row $(-y,x)$.
The kernel of the row is $S(x,y)^{\mathsf T}$ because $x,y$
are regular parameters in the domain $S$, and the image of
the column is $x S(x,y)^{\mathsf T}$. The cokernel of the row
is $k$. The dualized first map is injective, so the only nonzero
Ext groups are $S/(x)$ and $k$ in degrees one and two.
Shifting by $2$ in \ref{res:D05} gives the two claimed degrees
of $\mathcal D_R$. This matches \ref{res:D06} since $kx$ is simple,
and the dual of the top cohomology of $k[[y]]$ is $k[[y]]$
by \ref{res:D03} over that DVR and \ref{res:D04}.
\end{proof}
